% ============================================================================
%  Capitolo 21 -- Ingresso e uscita                              (Lezione 21)
% ============================================================================
\chapter{Ingresso e uscita}
\label{ch:io}

\epigraphline{Input! Ho bisogno di input!}{Numero 5, nel film
\emph{Corto circuito}, John Badham, 1986}

\begin{chapgoals}
Alla fine di questo capitolo saprete
\begin{itemize}
\item descrivere come si forma un'immagine su un CRT e su un LCD, e che cosa aggiunge una matrice attiva TFT;
\item spiegare il percorso della pressione di un tasto: matrice, codificatore, registro a scorrimento, interfaccia seriale;
\item dire che cosa invia un mouse al calcolatore, e quanto spesso;
\item descrivere le fasi della stampa laser;
\item spiegare come il kit \prg legge il suo joystick, e confrontare polling, interruzioni e DMA.
\end{itemize}
\end{chapgoals}

% ----------------------------------------------------------------------------
\section{Riscaldamento sulla macchina}

\asmfile[caption={Leggere il mondo, una volta per frame.},label=lst:io:warm]{code/l21/s04.S}

Leggete il joystick fin dai primi capitoli (Listato~\ref{lst:io:warm}). Ogni
pulsante è un interruttore su un piedino di ingresso del chip: premuto o
rilasciato, un bit. Il runtime legge i piedini, che sono registri mappati in
memoria configurati come ingressi, e li impacchetta in una parola: una
rappresentazione decodificata, nel senso del Capitolo~\ref{ch:combinational}.
E il gioco chiede una volta per frame: questo si chiama \term{polling}
(interrogazione ciclica).

Il listato ripete anche una lezione del Capitolo~\ref{ch:instructions}. La
maschera dei pulsanti è tenuta in \reg{s0}, un registro salvato, e quindi
sopravvive alla chiamata che suona la nota: il segnalino si muove con
\key{LEFT} anche mentre suona \key{A}. In \reg{t0} la maschera andrebbe persa
alla chiamata. (Il programma è stato compilato come cartuccia ed eseguito
sull'esecutore senza interfaccia di \prgqt.)

Una tastiera ha un centinaio di tasti, un mouse si muove in due dimensioni, un
display ha migliaia di pixel: come parlano con un processore?

% ----------------------------------------------------------------------------
\section{I display}

\subsection{Il video: un'introduzione}

I primi display mostravano soprattutto caratteri, in bianco e nero; oggi sono
grafici e a colori. Un display grafico è una matrice di punti, i \term{pixel},
ciascuno acceso con un'intensità diversa, e un'immagine si compone accendendo
e spegnendo i pixel. Una diversa densità di pixel dà una diversa risoluzione; i
colori vanno da 256 (8 bit) a 16 milioni (24 bit) e oltre. La dimensione si
misura in pollici lungo la diagonale. Il display del kit \prg ha una diagonale
di 2,8 pollici, 320 per 240 pixel e 65\,536 colori.

\subsection{Il tubo a raggi catodici}

\begin{table}[htb]
\centering\small
\caption{Breve storia del tubo a raggi catodici.}\label{tab:io:crt}
\begin{tabularx}{\textwidth}{@{}llL@{}}
\toprule
\textbf{Anno} & \textbf{Chi} & \textbf{Che cosa}\\
\midrule
primi anni 1870 & William Crookes & il tubo di Crookes: il primo tubo a raggi catodici, uno strumento di laboratorio\\
1897 & Karl Ferdinand Braun & il tubo di Braun e il primo oscilloscopio a raggi catodici: la prima applicazione pratica\\
1907 & Boris Rosing & un CRT mostra semplici forme geometriche da un segnale video ricevuto: la radice della televisione\\
1934 & Telefunken & il primo televisore con un CRT sul mercato, in Germania\\
anni 1970--1990 & l'industria dei calcolatori & il monitor CRT: il display di ogni terminale e personal computer\\
\bottomrule
\end{tabularx}
\end{table}

In un \term{tubo a raggi catodici} (CRT; Tabella~\ref{tab:io:crt},
Figura~\ref{fig:io:crt}) il catodo è un piccolo elemento metallico riscaldato
fino all'incandescenza, che emette elettroni per effetto termoionico. Dentro
il tubo gli elettroni sono diretti in un fascio da un'alta differenza di
potenziale tra catodo e anodo, e altri campi elettromagnetici mettono a fuoco
il fascio. Il fascio, il «pennello elettronico», è deflesso da campi magnetici
(la forza di Lorentz) o da campi elettrici, e colpisce un punto sulla
superficie interna dello schermo, che è rivestita di fosforo e si illumina.
Far scorrere il fascio riga per riga, molte volte al secondo, dipinge
l'immagine.

\begin{figure}[htb]
\centering
\begin{subfigure}[b]{0.3\textwidth}\centering\includegraphics[width=\textwidth,height=36mm,keepaspectratio]{mac_classic.jpg}\caption{Un Macintosh compatto: un CRT da 9 pollici}\end{subfigure}\hfill
\begin{subfigure}[b]{0.36\textwidth}\centering
\begin{tikzpicture}[x=0.42cm,y=0.42cm,font=\sffamily\tiny]
  \draw[prgsteel,thick,fill=prgpanel] (0,-0.9) -- (4.2,-0.9) -- (8,-2.9) -- (8.5,-2.9) arc (-35:35:5.05) -- (8,2.9) -- (4.2,0.9) -- (0,0.9) -- cycle;
  \draw[prgred,thick] (0.5,-0.25) rectangle (1.3,0.25);
  \fill[prgsteel] (2.0,1.1) rectangle (3.4,1.6) (2.0,-1.6) rectangle (3.4,-1.1);
  \fill[prgamber] (4.4,1.5) rectangle (5.6,2.0) (4.4,-2.0) rectangle (5.6,-1.5);
  \draw[prggreen,very thick,dashed] (1.3,0) -- (4.6,0) .. controls (6,0.2) .. (9.35,1.6);
  \fill[prggreen] (9.35,1.6) circle (2.5pt);
  \node[anchor=north] at (0.9,-1.0) {catodo};
  \node[anchor=south] at (2.7,1.6) {fuoco};
  \node[anchor=south] at (5.0,2.0) {deflessione};
  \node[anchor=west] at (9.6,1.6) {fosforo};
  \node[anchor=north] at (5.2,-3.0) {fascio di elettroni, anodo ad alta tensione};
\end{tikzpicture}
\caption{Lo schema di un CRT}\end{subfigure}\hfill
\begin{subfigure}[b]{0.3\textwidth}\centering\includegraphics[width=\textwidth,height=36mm,keepaspectratio]{crt_scope.jpg}\caption{Un oscilloscopio usato come display}\end{subfigure}
\caption{Il tubo a raggi catodici.}
\label{fig:io:crt}
\end{figure}

\subsection{Display a cristalli liquidi}

I monitor LCD hanno sostituito quasi del tutto il tubo a raggi catodici: danno
un'immagine migliore, con meno ingombro, peso e potenza, e affaticano meno gli
occhi. I \term{cristalli liquidi} sono molecole organiche viscose che scorrono
come un liquido ma hanno una struttura spaziale. Un campo elettrico cambia
l'orientamento delle molecole, e con esso una proprietà ottica del cristallo:
l'angolo di polarizzazione della luce che lo attraversa. Un cristallo è una
valvola per la luce, aperta e chiusa da una tensione.

\subsection{Dentro un pannello LCD}

Lo schermo è fatto di molte celle, spesso una per pixel, in cui sono
intrappolati i cristalli liquidi (Figura~\ref{fig:io:lcd}). Ogni cella ha
contatti elettrici, così che si può applicare un campo al liquido che
contiene. Le celle stanno tra due filtri polarizzatori i cui assi sono
perpendicolari tra loro, con una luce, la \term{retroilluminazione}, dietro
quello posteriore. La polarizzazione decide se la luce passa: quando i
cristalli non sono pilotati la luce viene ruotata e attraversa il polarizzatore
anteriore; quando sono pilotati viene completamente bloccata. In uno schermo a
colori ogni cella è divisa in tre sezioni con un filtro rosso, uno verde e uno
blu: i tre campi del colore RGB565 del Capitolo~\ref{ch:information} finiscono
qui, come tre tensioni.

\begin{figure}[htb]
\centering
\begin{subfigure}[b]{0.46\textwidth}\centering\includegraphics[width=\textwidth,height=38mm,keepaspectratio]{lcd_layers.jpg}\caption{Gli strati di un LCD}\end{subfigure}\hfill
\begin{subfigure}[b]{0.46\textwidth}\centering\includegraphics[width=\textwidth,height=38mm,keepaspectratio]{lcd_pixels.jpg}\caption{Sottopixel rossi, verdi e blu}\end{subfigure}
\caption{Display a cristalli liquidi.}
\label{fig:io:lcd}
\end{figure}

\subsection{Pilotare le celle: dai segmenti alla matrice attiva}

Nelle calcolatrici e negli orologi digitali c'è un contatto elettrico per
segmento, il che diventa impossibile al crescere del numero di elementi. Una
\emph{matrice passiva} ha un gruppo di contatti per riga e uno per colonna,
non per pixel; viene pilotato un solo pixel alla volta, e gli altri devono
ricordare il proprio stato. Le conseguenze sono la \emph{persistenza}, le scie
dietro gli oggetti in movimento, e la diafonia tra celle vicine.

Gli schermi ad alta risoluzione usano una \term{matrice attiva}. Il pannello
contiene una pellicola sottile di transistor (TFT, \emph{thin-film
transistor}) su uno strato semiconduttore trasparente; il campo è applicato
esattamente dove serve da un transistor in ogni cella, e il transistor
memorizza lo stato elettrico del suo pixel mentre gli altri vengono
aggiornati. Ogni pixel è, in effetti, una cella di memoria analogica. Il
display del kit è un pannello TFT con un controller ILI9341. Nei più recenti
display OLED ogni pixel emette la propria luce, e non serve alcuna
retroilluminazione.

% ----------------------------------------------------------------------------
\section{Tastiera e mouse}

\subsection{La tastiera}

La tastiera deriva da quella delle macchine da scrivere, ed è il principale
dispositivo per l'ingresso di caratteri. I suoi tasti sono raggruppati: tasti
alfanumerici; tasti speciali come invio e backspace; tasti freccia e tasti
funzione. Tastiera e video non sono collegati direttamente tra loro: è il
processore che riproduce sul video l'informazione battuta sulla tastiera.
L'eco di un tasto sullo schermo è un programma: leggere un codice, disegnare
un carattere. Le tastiere di oggi si collegano attraverso USB, o senza fili.

\subsection{Il percorso della pressione di un tasto}

\begin{figure}[htb]
\centering
\begin{tikzpicture}
  \foreach \k/\x/\y in {Q/0/0,W/1/0,E/2/0,R/3/0,A/0.3/1,S/1.3/1,D/2.3/1,F/3.3/1,Z/0.6/2,X/1.6/2,C/2.6/2,V/3.6/2}{
    \def\kd{D}\ifx\k\kd \def\c{prgamber}\else\def\c{prgsteel}\fi
    \node[draw=white,fill=\c,text=white,rounded corners=1pt,minimum size=5.5mm,font=\ttfamily\bfseries\small] at (\x*0.62,-\y*0.62) {\k};
  }
  \node[lbl,text=prgcomment] at (1.2,-1.9) {\emph{righe e colonne di interruttori}};
  \node[blkacc,text width=27mm,minimum height=14mm] (e) at (5.2,-0.6) {\textbf{codificatore}\\\footnotesize da riga e colonna a un codice di $n$ bit};
  \node[blkvio,text width=27mm,minimum height=14mm] (s) at (8.6,-0.6) {\textbf{registro a scorrimento}\\\footnotesize $n$ bit da parallelo a serie};
  \node[blkgrn,text width=27mm,minimum height=14mm] (u) at (12.0,-0.6) {\textbf{interfaccia seriale}\\\footnotesize p.es.\ un controller USB};
  \draw[->,thick] (2.9,-0.6) -- (e); \draw[->,thick] (e) -- (s); \draw[->,thick] (s) -- (u);
\end{tikzpicture}
\caption{Il percorso della pressione di un tasto.}
\label{fig:io:key}
\end{figure}

I tasti di una tastiera sono interruttori disposti in righe e colonne
(Figura~\ref{fig:io:key}). Un tasto premuto unisce una riga a una colonna; un
\emph{codificatore} trasforma quella combinazione di riga e colonna in un
codice di tasto di $n$ bit; un \emph{registro a scorrimento} porta gli $n$ bit
da parallelo a seriale; e un'\emph{interfaccia seriale}, per esempio un
controller USB, li invia al calcolatore.

\begin{prgconcept}
Ogni blocco della Figura~\ref{fig:io:key} è una macchina che avete già
progettato. È il frutto dei Capitoli da \ref{ch:combinational}
a~\ref{ch:registers}.
\end{prgconcept}

\subsection{Il mouse}

Il mouse fu inventato negli anni Sessanta da Douglas Engelbart (il primo
prototipo è del 1964, il brevetto fu depositato nel 1967 e la celebre
dimostrazione pubblica è del 1968): fu il primo dispositivo di puntamento
della storia. Era una scatola di legno con due ruote sotto e un filo che
usciva dal retro, come la coda di un topo. Le due ruote rilevavano i movimenti
lungo l'asse orizzontale e quello verticale, e l'elettronica li trasformava in
segnali elettrici. In un mouse \emph{meccanico} una pallina fa girare due
rulli, ciascuno dei quali aziona un contatore; in un mouse \emph{ottico} una
sorgente di luce e una minuscola telecamera misurano il movimento della
superficie sottostante.

\subsection{Che cosa invia un mouse}

Ogni volta che si è spostato di una distanza minima, circa un centesimo di
pollice, chiamata \emph{mickey}, il mouse invia tre byte: di quante unità si è
spostato lungo $x$ dall'ultimo rapporto; lo spostamento lungo $y$; e lo stato
dei pulsanti. I due spostamenti sono numeri con segno in complemento a due
(Capitolo~\ref{ch:integers}). I mouse moderni riferiscono fino a mille volte
al secondo. Dentro, il mouse è fatto di due contatori, un registro a
scorrimento che mette in serie i tre byte, e un'interfaccia seriale; dal lato
del calcolatore il software di controllo, il \term{driver di periferica},
sovrintende alla comunicazione con il processore.

% ----------------------------------------------------------------------------
\section{Le stampanti}

\subsection{La stampante}

Una stampante è un dispositivo di uscita che stampa informazione su carta. La
sua velocità si misura in righe al minuto, caratteri al secondo o pagine al
minuto; la sua risoluzione, cioè la qualità della stampa, in punti per
pollice. Le stampanti a margherita e a catena sono ormai in disuso; una
stampante \emph{ad aghi} colpisce con degli aghi un nastro inchiostrato; una
stampante \emph{a getto d'inchiostro} lancia minuscole gocce di inchiostro; e
poi c'è la stampante laser. Una stampante è anche un piccolo calcolatore: ha
processore, memoria e firmware propri.

\subsection{La stampante laser}

La stampante laser deriva dalla fotocopiatrice
(Figura~\ref{fig:io:laser}). Funziona in sei fasi.

\begin{enumerate}
\item Un tamburo viene caricato elettricamente, a circa 1\,000 volt.
\item Un raggio laser, riflesso da uno specchio rotante, scandisce la
  lunghezza del tamburo; il raggio è modulato in modo da produrre una
  configurazione di punti chiari e scuri.
\item I punti colpiti dal raggio perdono la carica. Dopo ogni riga il tamburo
  ruota di una frazione di grado.
\item Il toner è attratto dai punti ancora carichi: si forma un'immagine della
  riga.
\item Il tamburo viene premuto contro la carta e vi trasferisce il toner;
  rulli riscaldati lo fissano in modo permanente.
\item Il tamburo viene scaricato e pulito.
\end{enumerate}

La pagina viene prima costruita come bitmap nella memoria della stampante: un
frame buffer per la carta.

\begin{figure}[htb]
\centering
\includegraphics[width=0.5\textwidth]{laser_printer.jpg}
\caption{Il percorso della carta in una stampante laser.}
\label{fig:io:laser}
\end{figure}

% ----------------------------------------------------------------------------
\section{Ritorno alla macchina}

\subsection{I dispositivi del kit PRG32}

\begin{table}[htb]
\centering\small
\caption{I dispositivi del kit \prg.}\label{tab:io:kit}
\begin{tabularx}{\textwidth}{@{}P{0.2\textwidth}lLL@{}}
\toprule
\textbf{Dispositivo} & \textbf{Direzione} & \textbf{Come è collegato} & \textbf{Come lo usa una cartuccia}\\
\midrule
joystick e pulsanti & ingresso & un interruttore per piedino di ingresso (GPIO) & \code{prg32\_input\_read}: un bit per pulsante\\
display, TFT $320\times 240$ & uscita & bus seriale SPI, con DMA & le chiamate \code{prg32\_gfx\_*} disegnano nel frame buffer\\
amplificatore e altoparlante & uscita & bus seriale I2S & le chiamate \code{prg32\_audio\_*}: note, campioni, musica\\
radio Wi-Fi & ingresso e uscita & dentro il chip & caricamento delle cartucce, lo Store, multigiocatore\\
porta seriale USB & ingresso e uscita & USB & programmazione del firmware, messaggi di debug\\
\bottomrule
\end{tabularx}
\end{table}

Un dettaglio della Tabella~\ref{tab:io:kit} era stato promesso nel
Capitolo~\ref{ch:electronics}: un interruttore rimbalza per qualche
millisecondo quando viene premuto, e campionarlo una volta per frame, ogni
poche decine di millisecondi, nasconde il rimbalzo.

\subsection{Tre modi di trattare un dispositivo}

\begin{description}[style=nextline,leftmargin=1.5em]
\item[Polling (I/O programmato).] Il programma chiede al dispositivo, ancora e
  ancora, se ha qualcosa. È semplice, ed è ciò che fa un gioco una volta per
  frame; è uno spreco quando non succede niente.
\item[Interruzioni.] Il dispositivo chiede attenzione quando gli serve; il
  processore sospende ciò che sta facendo, serve il dispositivo e riprende
  (Capitolo~\ref{ch:instructions}).
\item[Accesso diretto alla memoria.] Il dispositivo scambia da solo interi
  blocchi di dati con la memoria; il processore avvia soltanto il
  trasferimento e viene avvisato quando termina (Capitolo~\ref{ch:memory}).
\end{description}

\begin{prgmachine}
Sul kit tutti e tre sono al lavoro nello stesso momento: il gioco interroga il
joystick, il timer che scandisce i frame interrompe, e il display e il suono
sono alimentati dal DMA.
\end{prgmachine}

% ----------------------------------------------------------------------------
\section{Verificate di aver capito}

\begin{quiz}{Che cosa fa il codificatore di un controller di tastiera?}
\item Illumina il tasto premuto
\item Trasforma una combinazione di riga e colonna in un codice di tasto
\item Invia il codice al monitor
\item Conta per quanto tempo un tasto resta premuto
\end{quiz}

\begin{quiz}{In un LCD a colori, un pixel è fatto di:}
\item una cella che cambia colore
\item tre celle con filtri rosso, verde e blu
\item un punto di fosforo colpito da elettroni
\item un transistor che emette luce
\end{quiz}

\begin{quiz}{Perché un LCD a matrice passiva mostra scie dietro gli oggetti in movimento?}
\item La retroilluminazione è troppo debole
\item I pixel sono pilotati uno alla volta e nel frattempo devono mantenere il proprio stato
\item Ha troppo pochi colori
\item La scheda video è lenta
\end{quiz}

\begin{quiz}{Un gioco \prg legge il joystick una volta per frame. Questa tecnica si chiama:}
\item accesso diretto alla memoria
\item ingresso a interruzioni
\item polling
\item mappatura in memoria
\end{quiz}

\section{Esercizi}

\begin{exercise}[prima su carta, poi su \prg]
\begin{enumerate}[nosep]
\item Una tastiera ha 104 tasti: quanti bit richiede il codice di tasto?
  Quante righe e colonne potrebbe avere la matrice?
\item Una pagina di 8 per 11 pollici a 600 punti per pollice, un bit per
  punto: quanta memoria serve alla stampante per la bitmap?
\item Muovete un puntatore sullo schermo di \prg con il joystick, come farebbe
  un mouse: tenete $x$ e $y$, e limitateli allo schermo.
\item Fate agire il pulsante \key{A} una sola volta per pressione (rilevamento
  del fronte, Capitolo~\ref{ch:sequential}).
\item Disegnate un punteggio a sette segmenti con dei rettangoli, usando la
  tabella del transcodificatore del Capitolo~\ref{ch:combinational}.
\end{enumerate}
\expected{Su carta: 7 bit e, per esempio, 8 righe per 13 colonne;
$4\,800\times 6\,600$ punti, cioè 31\,680\,000 bit, circa 4~MB. Su \prg: un
puntatore che resta sullo schermo, e un pulsante che scatta una volta per
pressione.}
\end{exercise}

\section{Glossario}

\begin{glossarytable}
cathode-ray tube & tubo a raggi catodici\\
liquid crystals & cristalli liquidi\\
backlight & retroilluminazione\\
active matrix & matrice attiva\\
keyboard / key & tastiera / tasto\\
pointing device & periferica di puntamento\\
dot-matrix / ink-jet printer & stampante ad aghi / a getto d'inchiostro\\
drum / toner & tamburo / toner\\
polling & interrogazione ciclica\\
device driver & driver di periferica\\
\end{glossarytable}

\section{Conclusioni e letture}

I dispositivi di ingresso/uscita sono un componente fondamentale dei
calcolatori. La tecnologia dei display ha fatto progressi enormi, dal fascio
di elettroni a un transistor dietro ogni pixel. Tastiera, mouse e stampante
sono cambiati poco nella forma, e dentro sono costruiti con codificatori,
contatori e registri a scorrimento. Un dispositivo raggiunge il processore
attraverso un controller e un bus, ed è trattato con il polling, con le
interruzioni o con il DMA. Sul kit \prg l'intera catena, da un interruttore a
un pixel, può essere seguita istruzione per istruzione.

\paragraph{Letture.} \citet{tanenbaum2013}, paragrafo~2.4 (ingresso/uscita);
\citet{stallings2021}, il capitolo sull'ingresso/uscita; \emph{Playing with
RISC-V}, capitolo~6, «State and Input» \citep{prg32textbook}.

Il quarto blocco è completo. Il prossimo capitolo conclude \vermone, e quello
successivo pubblica il gioco.
